\documentclass[a4paper]{scrreport}

\usepackage[T1]{fontenc}
\usepackage[utf8]{inputenc}
\usepackage[nswissgerman]{babel}
\usepackage[babel,german=swiss]{csquotes}
\usepackage[onehalfspacing]{setspace}
\usepackage[style=verbose-ibid]{biblatex} % -ibid


% Literaturverzeichnis hinzufügen
\addbibresource{bibliography.bib}


% Einstellungen für Darstellung der Fussnoten
\usepackage[bottom,norule,hang,perpage]{footmisc}
\setlength{\footnotemargin}{2mm} % Spacing in Fussnoten zwischen der Zahl und dem Text
% \renewcommand{\footnoterule}{\rule{45mm}{0.1mm}{\vspace*{5mm}}} % Trennlinie über Footer


% Einstellunge für Zitierart
\renewcommand*{\labelnamepunct}{\addcolon\addspace} % Interpunktion in Zitaten zwischen Autor und Titel
\renewcommand{\newunitpunct}{\addcomma\addspace} % Zitatelemente mit einem Komma trennen (nicht der standardmässig Punkt)
\renewcommand{\subtitlepunct}{\addperiod\addspace} % Interpunktion in Zitaten zwischen Titel und Untertitel


% Spacing des Papiers
\usepackage{geometry}
\geometry{a4paper,top=30mm,right=27mm,bottom=30mm,left=27mm}


% Überschriftentiefe die im Inhaltsverzeichnis nummeriert werden
\setcounter{secnumdepth}{3}


% Überschriftentiefe die im Inhaltsverzeichnis angezeigt werden
\setcounter{tocdepth}{3}


% Abbildungen (figure) innererhalb einer Section durchnummerieren
% https://tex.stackexchange.com/a/425603
% \let\counterwithout\relax
% \let\counterwithin\relax
% \usepackage{chngcntr}
% \counterwithin{figure}{section}
% \counterwithin{table}{section}


% Sections durchgehend nummerieren und nicht von Chapters unterbrechen
% \usepackage{chngcntr} % chngcntr bereits eingebunden
% \counterwithout{section}{chapter}


% Package um Links im PDF-Dokument zu verwenden
\usepackage{url}


% Package um Bilder (PNG) und PDFs als figures einzubinden
\usepackage{graphicx}


% Placeholder for figures: \includegraphics[width=.5\textwidth]{example-image-a}
% https://tex.stackexchange.com/a/413186
\usepackage{mwe}


% Package um Symbole für Musik bereitzustellen: \flat \sharp
% \usepackage{textcomp}


% Package um Symbole Noten bereitzustellen: \quarternote \twonotes
% \usepackage{wasysym}


% Package um Generalbassbezifferungen im Fließtext besser zu verwenden
% https://tex.stackexchange.com/a/267644
% https://tex.stackexchange.com/a/267618
% \usepackage{amsmath} % \overset{above}{main}


% Command für inline Code im Fließtext
\newcommand{\code}{\texttt}


% Package und Command um Text einzukreisen für Bassstufen
% https://tex.stackexchange.com/a/199056
\usepackage{tikz}
\newcommand*\circled[1]{\tikz[baseline=(char.base)]{\node[anchor=text, shape=circle,draw, inner sep=0pt, minimum size=1.2em] (char) {#1\strut};}}


% Git Hash der Version erkennen und auf \gitAbbrevHash legen
\usepackage{xstring}
\usepackage{catchfile}
\CatchFileDef{\gitHash}{.git/refs/heads/master}{}
\newcommand{\gitAbbrevHash}{\StrLeft{\gitHash}{7}}


% Alle Quellen im Literaturverzeichnis anzeigen, auch nicht verwendete
\nocite{*}


% Einstellungen für das Stichwortverzeichnis
% \usepackage{makeidx}
% \makeindex
% \renewcommand{\indexname}{Stichwortverzeichnis}


% Nummerierung von Sections ohne die Zählung von Chapter
\renewcommand*\thesection{\arabic{section}}


% Spacings in Tabellen
\setlength{\tabcolsep}{1em} % horizontal padding
{\renewcommand{\arraystretch}{1.5}% vertical padding


% Fussnoten in Tabellen aktivieren mit dem Command \tablefootnote
\usepackage{tablefootnote}


% Zeilenumbrüche in Tabellen erlauben mit \thead{Mixo-\\lydisch} und |makecell{}
% https://tex.stackexchange.com/a/176780
\usepackage{makecell}


% Fügt Fraktur Schriftarten hinzu: \textfrak{Fraktur} \textswab{Schwabacher} \textgoth{Gothic}
% \usepackage{yfonts}


% "Abbildung ..." unter den figures umbenennen
\usepackage[figurename=Notenbeispiel]{caption}
\renewcaptionname{nswissgerman}{\listfigurename}{Verzeichnis der Notenbeispiele}


% Gestrichelte horizontale Linien in Tabellen mit \hdashline
\usepackage{arydshln}


% Abstände im Inhaltsverzeichnis zwischen Zahl und Titel anpassen
\usepackage{tocloft}
\setlength{\cftsecnumwidth}{1.5em}
\setlength{\cftsubsecnumwidth}{2.1em}
\setlength{\cftsubsubsecnumwidth}{2.75em}


% Sicherstellen, dass figures in der section ausgegeben werden zu der sie gehören
\usepackage[section]{placeins}




\title{Unbekannter Titel}
\subtitle{Master Musikpädagogik}
\subject{Masterthesis}
\titlehead{
	\hfill
	\includegraphics{img/Logo_Musikhochschule.pdf}
}
\author{
	Katharina Jäckle\\
	Hochschule für Musik Freiburg
}
\dedication{
	{
		\Huge
		Nochmal der Titel hier\\
	}
	\vspace{0.75cm}
	{
		\Large
		Und hier dann ein etwas längerer\\
		Untertitel der auch über mehrere\\
		Zeilen gehen darf\\
	}
	\vspace{1.5cm}
	{
		\normalsize
		Ausserdem noch eine kleine Danksagung \\
		oder weitere kurze Unteruntertitel hier\\
	}
	\vspace{12cm}
	{
		\footnotesize
		Text- und Notensatz mit \LaTeX, LilyPond und git.\\
		\small{Version: \gitAbbrevHash}
	}
}
\publishers{
	Gutachter: Andreas Doerne\\
}
\date{\today}



\begin{document}

	\maketitle

	\tableofcontents

    \addchap{Vorwort}

Danke an die Einflüsse, Anstöße und den Input meiner Professoren Wolfgang
Lessing und Andreas Doerne, die vorallem im musikpädagogischen Bereich prägend
waren aber sich auch auf mein künstlerisches Dasein ausgewirkt haben. 
An Andreas Doerne, der unsere Masterarbeit betreut
Wolfgang Lessing, der uns zum Wettbewerb gebracht hat etc.
Michael Stecher?
Und an meinen Kommilitonen Timo Langpap, mit dem ich Co-Working betrieben, mich
immer wieder ausgetauscht und diskutiert habe. Der außerdem mit seiner Biografie
noch einmal einen anderen Blick auf die Dinge hat und auch in seinem Genre dem
Jazz und Pop auch andere Gepflogenheiten gewohnt ist, was immer zu einem
spannenden Austausch geführt hat. 

Außerdem bedanke cih mich fürs Korrekturlesen und die Unterstützung in der
Nutzung des Programms LaTex bei Wolfgang Drescher, der mir bei all diesen Fragen
stets zur Seite stand. 

Durch unsere Gedanken zu einem eigenen Musikschulmodell, habe ich mir auch über
das gemeinsame Musizieren und Musiklernen in der Gemeinschaft viele Gedanken
gemacht.




\vspace{0.5cm}

\begin{flushright}
	{
		\small
		Freiburg im Breisgau, September 2021\\
		--- Katharina Jäckle
	}
\end{flushright}

	\addchap{Einleitung}
Das gemeinsame Musizieren zieht sich bei meinem Lebenslauf wie ein roter Faden
durch die ganze Laufbahn. Als Instrumentallehrerin habe ich allerdings auch
beobachtet, wie es für Schüler ist, die lediglich den
Einzel-Instruemntalunterricht kennen. Dabei ist mir aufgefallen, dass ein
wichtiger Bestandteil der für mich die Musik ausmacht fehlt. Nämlich das
gemeinsame Musizieren. Wenn man alleine den Aspekt des
Differenziellen Lernens betrachtet wird schnell klar, das Einzelunterricht
alleine nicht ausreicht um einen Schüler motiviert am Instrument zu halten und
ihn weiter zu bringen. Da sich auf der sozialen Ebene im Miteinander so viel
ereignet, kann eine einzelne Lehrkraft nicht alle Ebenen abdecken und bedienen.
Oft ist in unserem Musikschulsystem das gemeinsame Musizieren aber gar nicht
vorgesehen, sondern vielmehr ein optionales Zusatzangebot, im Orchester, im Chor
oder in der Band zu spielen. Dieses Zusatzangebot ist auch nicht immer gegeben,
je nach Größe der Musikschule und schon sind die Schüler schnell auf sich und
ihren Instrumentallehrer allein gestellt. In anderen Ländern gibt es andere
Sitten und somit auch verschiedene Systeme, in denen Instrumentalschüler mehr im
Miteinander das Instrument lernen. Einige will ich deshalb näher beleuchten und
unteruchsen, inwieweit sie für die Lernenden wertvoll sind. Im Vordergrund
meiner Arbeit soll aber stehen, warum das gemeinsame Musizieren so
wichtig ist und das werde ich ich im Folgenden beleuchten. 


"(...) dass Muisk in ihrem Wesen Mitteilung, Kommunikation, Gestaltung, das
Vermitteln einer Botschaft ist, die in Sternstunden über das Erleben der
Einzelnen hinausreicht und sie miteinander verbinden kann." S. 11

"Musik als Kommunikationsmittel" S. 17

Aristoteles: "der Schlüssel zu einem glücklichen Leben besteht im Miteinander,
das Leben in der Gemeinschaft unter Freunden." (S. 9 "Die Kunst zu unterrichten)

"Sie muss gelernt und ein Leben lang geübt werden." (auch S. 9)


"Musik ist Teil der Symbolwelt des Menschen: Mitteilung, Kommunikation,
Interaktion." Wolfgang Suppan, (S. 91 Andreas Doerne)



	\clearpage
	\addcontentsline{toc}{chapter}{Literatur}
	\printbibliography

	% \clearpage
	% \addcontentsline{toc}{chapter}{Stichwortverzeichnis}
	% \printindex

	% \clearpage
	% \addcontentsline{toc}{chapter}{Verzeichnis der Notenbeispiele}
	% \listoffigures

	\section*{Erklärung}

\pagenumbering{gobble}

Hiermit erkläre ich, dass ich diese Arbeit selbständig verfasst habe, keine
anderen als die angegebenen Quellen und Hilfsmittel verwendet habe und alle
Stellen, die wörtlich oder sinngemäß aus veröffentlichten Schriften entnommen
wurden, als solche kenntlich gemacht habe. Darüber hinaus erkläre ich, dass
diese Arbeit nicht, auch nicht auszugsweise, bereits für eine andere Prüfung
angefertigt wurde.


\end{document}
